\documentclass[a4paper,11pt]{article}
\usepackage{iftex}
\ifPDFTeX
  \usepackage[utf8]{inputenc}
  \usepackage[T1]{fontenc}
  \usepackage{lmodern}
\else
  \usepackage{fontspec}
  \setmainfont{lmroman10-regular.otf}[
    BoldFont=lmroman10-bold.otf,
    ItalicFont=lmroman10-italic.otf,
    BoldItalicFont=lmroman10-bolditalic.otf]
  \setmonofont{lmmono10-regular.otf}
\fi
\usepackage[croatian]{babel}
\usepackage[margin=24mm]{geometry}
\usepackage{enumitem}
\usepackage[hidelinks]{hyperref}
\setlength{\parindent}{0pt}
\setlength{\parskip}{6pt}
\setlist[itemize]{nosep,leftmargin=16pt,topsep=3pt}
\setlist[enumerate]{nosep,leftmargin=18pt,topsep=3pt}
\hypersetup{pdftitle={Data Monitoring Platform - Izvješće 1},pdfauthor={Leonard Lečić i Dario Golubović}}
\AtEndDocument{\clearpage\ifnum\value{page}=3\relax\else
\PackageError{KT1}{Izvjesce mora imati dvije stranice}{Prilagodite duljinu sadrzaja.}\fi}
\begin{document}
\begin{center}
{\Large\bfseries Data Monitoring Platform}\\[5pt]
{\large Izvješće 1 -- Prva kontrolna točka}\\[7pt]
Leonard Lečić i Dario Golubović\\
Skriptni jezici \quad | \quad 28. listopada 2026.
\end{center}

\section*{1. Opis problema i cilj}
Projekt je sustav Python skripti za automatsko prikupljanje, obradu i
praćenje podataka iz više vanjskih API-ja. Cilj je objediniti podatke,
izgraditi povijest, izračunavati promjene i upozoravati na zadane uvjete.
Web sučelje služit će pregledu rezultata i konfiguraciji pravila.

Naglasak je na skriptiranju i automatizaciji: raspoređenom izvršavanju,
obradi JSON i CSV zapisa, normalizaciji, statistici, generiranju sažetaka
i automatskim obavijestima. Dohvat i obrada radit će neovisno o tome je li
web sučelje otvoreno. Rješenje se razvija samostalno, s prilagodljivim
modulima za različite vrste podataka.

\section*{2. Opseg i automatizacija}
Planirana su tri API sučelja: \textbf{cijene.dev} za cijene proizvoda te
\textbf{Open-Meteo Weather} i \textbf{Open-Meteo Air Quality} za vrijeme
i kvalitetu zraka. Pratit će se odabrani proizvodi i poslovnice te
vremenski podaci i PM2.5/PM10 za Rijeku i Zagreb. Cijene će se dohvaćati
dnevno, a ostali podaci jednom na sat, uz prilagodljiv raspored.

Predviđene su sljedeće mogućnosti:
\begin{itemize}
\item konfiguracija izvora, lokacija, intervala i pragova bez izmjene koda;
\item automatski i ručni dohvat, ponovni uvoz te obrada spremljenih podataka;
\item povijesna pohrana, provjera valjanosti i uklanjanje duplikata;
\item osnovna statistika, klizni prosjeci, trendovi i otkrivanje odstupanja;
\item pravila upozorenja i automatsko slanje e-pošte preko SMTP-a;
\item dnevni i tjedni sažeci te uvoz i izvoz u CSV/JSON formatu;
\item evidencija izvršavanja, provjera svježine podataka i sigurnosne kopije.
\end{itemize}

Sučelje će prikazivati grafove i upozorenja, omogućiti pretragu i filtriranje
te povremeno osvježavati prikaz bez ponovnog učitavanja stranice. Ručni
zahtjev za novim podacima pokretat će isti postupak kao raspoređeni posao.

\section*{3. Tehnologije}
Glavni jezik bit će \textbf{Python 3.12+}. Koristit će se \textbf{httpx} za
HTTP zahtjeve, \mbox{\textbf{Pydantic}} za validaciju, \textbf{SQLite} i
\texttt{sqlite3} za pohranu, \textbf{pandas} za analizu te
\textbf{APScheduler} za raspored. Standardni moduli \texttt{argparse},
\texttt{tomllib}, \texttt{logging}, \texttt{csv}, \texttt{json} i
\texttt{smtplib} pokrivat će naredbe, konfiguraciju, logove, izvoz i e-poštu.

\textbf{PowerShell} skripte pojednostavit će pokretanje i provjere.
\textbf{FastAPI} i \textbf{Uvicorn} izlagat će rezultate, a
\mbox{\textbf{JavaScript}} i \textbf{Chart.js} prikazivati ih uz HTML/CSS.
Testovi će se pisati u \textbf{pytestu}, a kod verzionirati Gitom.

\newpage
\section*{4. Arhitektura i obrada}
Sustav će biti podijeljen na module za dohvat, normalizaciju, pohranu,
analizu, pravila, obavijesti i API. Zaseban Python proces izvršavat će
raspoređene poslove. Naredbe za dohvat, analizu, izvoz i sigurnosnu kopiju
koristit će iste funkcije kao automatizirani poslovi. Novi izvor dodavat
će se zasebnim adapterom koji vraća zajednički format zapisa.
\begin{center}
\small Dohvat $\rightarrow$ validacija $\rightarrow$ pohrana
$\rightarrow$ analiza $\rightarrow$ pravila i obavijesti
\end{center}

Zapis će sadržavati izvor, identitet entiteta, metriku, vrijednost, jedinicu
te vrijeme podatka i preuzimanja. Baza će sadržavati i izvršavanja poslova,
pravila i događaje upozorenja. Cijene će se obrađivati tipom
\texttt{Decimal}, a vremena pohranjivati u UTC-u. Analiza će uspoređivati
iste metrike i jedinice; cijene istog proizvoda i pakiranja u istoj
poslovnici. Ponovni uvoz neće stvarati duplikate.

Pravila će obuhvaćati prelazak praga, značajnu postotnu promjenu i
zastarjele podatke. Evaluirat će se nakon dohvata, uz ograničavanje
ponavljanja obavijesti. Poslovi će imati timeout, ograničene ponovne
pokušaje i zapis ishoda. Pogreška jednog izvora neće zaustaviti ostale;
neuspjela e-pošta evidentirat će se za naknadno slanje.

\section*{5. Napredak u prvoj kontrolnoj točki}
Postavljen je repozitorij \texttt{data-monitoring-platform}. Definirani su
problem, opseg, tehnologije, struktura modula i plan razvoja za dva člana.
Razrađen je tok od dohvata do obavijesti te konfiguriranje i samostalno
pokretanje skripti. Implementacija integracija i aplikacijskih
funkcionalnosti slijedi u idućim kontrolnim točkama.

\section*{6. Plan i organizacija rada}
\begin{itemize}
\item \textbf{KT1 -- 1. studenoga:} repozitorij, definicija projekta,
arhitektura, izbor tehnologija i raspored rada.
\item \textbf{KT2 -- 27. studenoga:} API adapteri, konfiguracija, naredbe
za ručni dohvat, SQLite pohrana, raspoređivanje i logovi; testovi
neispravnih odgovora i ponovnog uvoza te početak prikupljanja povijesti.
\item \textbf{KT3 -- 18. prosinca:} analitika, pravila, SMTP obavijesti,
CSV/JSON izvoz i sažeci; API i sučelje s filtrima, grafovima i
automatskim osvježavanjem.
\item \textbf{KT4 -- 22. siječnja:} provjera svježine, sigurnosne kopije
i oporavak, integracijski testovi, ispravci i dokumentacija;
demonstracija cijelog automatiziranog postupka i rada pri greškama.
\end{itemize}

Rad će se podijeliti na dohvat, pohranu i raspoređivanje te na analitiku,
pravila, obavijesti i prikaz. Oba člana sudjelovat će u Python dijelu;
integracija, testiranje i dokumentacija provodit će se zajednički.
\end{document}
